\documentclass[10pt, aspectratio=1610]{beamer}
%\documentclass{article}
%\usepackage{beamerarticle}
\mode<article>
{
	\usepackage[a4paper]{geometry}
}



\edef\tlCourse{Промышленное программирование}
\edef\tlPartName{Модуль 1. Разработка модульных программ}
\edef\tlPartNumber{1}
\edef\tlThemeName{Системы сборки проектов на языке С/С++}
\edef\tlThemeNumberLocal{1}
\edef\tlThemeNumberGlobal{1}


%\documentclass[10pt, handout]{beamer}
%\setbeameroption{show notes}

%\documentclass[10pt, a4paper]{article}
%\usepackage{beamerarticle}



%Для XeLatex/+
%\usepackage{polyglossia}
%\setdefaultlanguage{russian}
%\setotherlanguage{english}
%%\setkeys{russian}{babelshorthands=true} 

\usepackage[russian]{babel}



\usepackage{fontspec}

\setmainfont{Times New Roman} 
\setsansfont{Arial} 
\setmonofont{Consolas}  



\mode<article>{
	
	\usepackage{hyperref}
	
}
\mode<presentation>{
	
	%\usetheme{Antibes}
	\usetheme{Berlin}
	
	\usefonttheme{professionalfonts} 
	\usefonttheme{serif} % default family is serif
	
	\usepackage{fontspec}
	\setmonofont{Cascadia Code SemiLight} %CMU Typewriter Text
	\setromanfont{DejaVu Serif}
	\setsansfont{DejaVu Sans}
	%\usecolortheme{spruce} %зеленая, плохой цвет в заголовках 
	%\usecolortheme{albatross} %синяя, пхоло виден черный цвет
	
	\setbeamercovered{dynamic}
	
}

\newlength{\parskipbackup}
\setlength{\parskipbackup}{\parskip}
\newlength{\parindentbackup}
\setlength{\parindentbackup}{\parindent}

\let\notebackup\note
\renewcommand{\note}[1]{\notebackup{%
		\setlength{\parindent}{3ex}%
		\setlength{\parskip}{0.5ex}%
		{\small{}#1}%
		\setlength{\parskip}{\parskipbackup}%
		\setlength{\parindent}{\parindentbackup}%
	}%
}

\mode<presentation>%
{
	
	%\setbeameroption{second mode text on second screen=right}
	%\setbeameroption{show notes on second screen}
	%\setbeameroption{show notes}
	\setbeamertemplate{note page}{
		
		%\vspace{1em}\insertframetitle
		
		%		\begin{wrapfigure}[12]{l}[0cm]{0.5\textwidth}
			%			\fbox{ \resizebox{0.45\textwidth}{!}{\insertslideintonotes{1}} }
			%		\end{wrapfigure}		
		\insertnote}
	
	
}


\newcommand{\MP}[1]{\mode<presentation>{#1} }
\newcommand{\MA}[1]{\mode<article>{#1} }

\newcommand{\ABS}[1]{\left| #1 \right|}
%\newcommand{\ABS}[1]{\mid #1 \mid}

\newcommand{\HREF}[2]{{\color{blue}\underline{\href{#1}{#2}}}}

\setbeamertemplate{caption}[numbered]


\usepackage{array}
\usepackage{tabularx}

\usepackage{verbatim}

\usepackage{ifthen}



\usepackage{tikz}
\usetikzlibrary{arrows.meta}
\usetikzlibrary{calc}
\usetikzlibrary{decorations}
\usetikzlibrary{decorations.pathreplacing}
\usetikzlibrary{matrix}
\usetikzlibrary{fit}

\newcommand{\rememb}[1]{\parbox{0pt}{}\tikz[remember picture,baseline=-0.5ex]{\draw node[inner sep=0pt, outer sep=0pt] (#1){\strut};}}



\usepackage{tikz-3dplot}
\usepackage{environ}
\usepackage{animate}





\usepackage{xcolor}


\usepackage{amsmath} %математические формулы
%\usepackage{amsfonts} %математические шрифты


\usepackage[e]{esvect}  %Красивая стрелочка вектора
%\let\oldvv\vv
\newcommand{\VV}[1]{\vv{#1\mathstrut}}

\usepackage{graphicx} %работа с каритнками




\usepackage{array}
\usepackage{multirow}



\usepackage{caption}
\DeclareCaptionLabelSeparator{dot}{~---~}            %Разделитель номер рисунка
%\captionsetup[figure]{justification=centering,labelsep=dot, format=plain}                        %Подпись рис. центр
%\captionsetup[table]{justification=raggedleft,labelsep=dot, format=plain, singlelinecheck=false} %Подпись табл. слева
\captionsetup[lstlisting]{justification=raggedleft,labelsep=dot, format=plain, font=normal, singlelinecheck=false}                     %Подпись рис. центр

\usepackage{indentfirst} %отступ первой строки


\usepackage[svgnames]{xcolor}


\usepackage{hyperref}



\def\sectionname{Раздел}
\def\subsectionname{Подраздел}


\NewDocumentCommand\TC{m+m+m}{%	
	\begin{columns}
		\begin{column}{#1\textwidth}
			#2
		\end{column}
		\begin{column}{\fpeval{1-#1}\textwidth}
			#3
		\end{column}
	\end{columns}
}

\NewDocumentCommand\TCT{m+m+m}{%	
	\begin{columns}[T]
		\begin{column}{#1\textwidth}
			#2
		\end{column}
		\begin{column}{\fpeval{1-#1}\textwidth}
			#3
		\end{column}
	\end{columns}
}


\usepackage{newfile}

\edef\LectionNumber{0}
\edef\LectionTheme{0}

\let\oldsection\section
\let\oldsubsection\subsection


\AtBeginDocument
{
	\newoutputstream{CONTENT}
	\openoutputfile{\jobname.gvr}{CONTENT}
	
	%\expandafter\addtostream{CONTENT}{\noindent\textbf{\Large Лекция \LectionNumber~---~\LectionTheme}\unexpanded{\setcounter{SEC}{0}}\par}
	\addtostream{CONTENT}{\protect\noindent%
						  	\protect\setcounter{SEC}{0}%
							{%
							\protect\large% 
							\protect\bfseries{}%
							Лекция \LectionNumber~---~\LectionTheme% 
							}% 
							\protect\par}
}

\renewcommand{\section}[2][]{
	\ifstrempty{#1}%
	{\oldsection{#2}}%
	{\oldsection[#1]{#2}}	
	\addtostream{CONTENT}{%
			\protect\noindent%
			\protect\stepcounter{SEC}%
			\protect\setcounter{SUB}{0}%
			\protect\hspace{3ex}%
			\protect\theSEC ~ #1 \protect\par%
			}
}

\renewcommand{\subsection}[1]{
	\oldsubsection{#1}
	\addtostream{CONTENT}{%
			\protect\noindent%
			\protect\stepcounter{SUB}%
			\protect\hspace{6ex}%
			\protect\theSEC{}.\protect\theSUB ~ #1 \protect\par%
			}
}


\newcommand{\Strut}{{\Large\strut}}

\newcommand\scb[1]{\left( #1 \right)}

\newcommand{\LINK}[2]{%
	\qrcode[height=1cm]{#1}\  \HREF{#1}{\parbox{0.8\textwidth}{#2}} \\[0.5em]
}

\NewDocumentCommand{\lecdni}{}{\LectionNumber}
\author{Гаврилов Андрей Геннадьевич}
\newcommand{\regals}{кандидат технических наук, доцент}
\institute{Кафедра Информационных технологий и вычислительных систем \\МГТУ~<<СТАНКИН>>}
\lecture{История компьютерной графики}{kghistory}\subtitle{Компьютерная графика}


\makeatletter
\newcommand*{\overlaynumber}{\number\beamer@slideinframe}
\makeatother





\usepackage{cprotect}

\usepackage{qrcode}

\newcommand{\QRFRAME}{%
    \begin{frame}[plain, noframenumbering]    	
	
	\centering
	Трансляция презентации (во время очных лекций)    
	
	~
	
	{\Large \ttfamily  https://clck.ru/3D3Efj  }
	
	~
	
	\tikz\node[inner sep=0pt,rounded corners=5mm, clip]{\qrcode[height=0.4\textwidth]{https://clck.ru/3D3Efj}}; 
	
	~	
	{\small
	При просмотре презентации в PDF для отображения анимаций на слайдах необходимо использовать Acrobat Reader, KDE Okular, PDF-XChange, Foxit Reader, браузер Firefox. Для браузеров на движке Chrome (Edge, Яндекс, Opera,~\dots) необходимо использовать \HREF{https://chromewebstore.google.com/detail/pdf-viewer/oemmndcbldboiebfnladdacbdfmadadm?hl=ru&utm_source=ext_sidebar}{PDF.js} c опцией <<Enable active content (JavaScript) in PDFs>>. }
	
	\end{frame}%
}

\newcommand{\IG}[2][1]{\includegraphics[width=#1\textwidth]{#2}}


%===========Настройка листингов==================

\usepackage{fancyvrb}
\usepackage{minted}

\AtBeginEnvironment{minted}{%
	\renewcommand{\fcolorbox}[4][]{#4}}

\mode<presentation>
{
	\newminted[cppcode]{c++}{linenos, autogobble=true, style=vs,tabsize=2,fontsize=\small,escapeinside=\@\@,xleftmargin=0em, numbersep=0.5em}

	
	\newminted[cppresumecode]{c++}{linenos, firstnumber=last,  autogobble=true, style=vs,tabsize=2,fontsize=\small,escapeinside=\@\@,xleftmargin=0em, numbersep=0.5em}
	
	\newminted[cmakecode]{cmake}{linenos, autogobble=true, style=default,tabsize=2,fontsize=\footnotesize,escapeinside=\@\@,xleftmargin=0em, numbersep=0.5em}
	
	\newminted[cmdcode]{Batch}{linenos, autogobble=true, bgcolor=red,  style=lightbulb,tabsize=2,fontsize=\small,escapeinside=\@\@,xleftmargin=0em, numbersep=0.5em}
}

\mode<!presentation>
{
	\newminted[cppcode]{c++}{linenos,  frame=lines, style=vs,tabsize=3,fontsize=\small,escapeinside=\@\@,xleftmargin=2em}
	\newminted[cppresumecode]{c++}{linenos, firstnumber=last, frame=lines, style=vs,tabsize=3,fontsize=\small,escapeinside=\@\@,xleftmargin=2em}
	
	\newminted[cmakecode]{cmake}{linenos, autogobble=true, style=vs,tabsize=2,fontsize=\small,escapeinside=\@\@,xleftmargin=2em}
	
	\newminted[cmdcode]{batch}{linenos, autogobble=true, style=lightbulb,tabsize=2,fontsize=\small,escapeinside=\@\@,xleftmargin=0em, numbersep=0.5em}
}

\newenvironment{codeblock}[2][]
{
	\VerbatimEnvironment
	\begin{block}{#1\hfill{}#2}
		\begin{cppcode}% 
}
{		
	\end{cppcode}
	\end{block}
}

\newenvironment{codeblockresume}[2][]
{
	\VerbatimEnvironment
	\begin{block}{#1\hfill{}#2}
		\begin{cppresumecode}% 
}
{		
		\end{cppresumecode}
		\end{block}
}


%===================================================




\graphicspath{{Images/}{Images/\jobname/}}

\date{\today}

\renewcommand{\LectionNumber}{\tlThemeNumberGlobal}
\renewcommand{\LectionTheme}{Тема \tlThemeNumberGlobal \\ \tlThemeName}
\title{\LectionTheme}
\subtitle{\tlCourse}


\usepackage{fancyvrb}
\usepackage{minted}


\usetikzlibrary {shapes.geometric}
\usetikzlibrary{shapes.misc}
\usetikzlibrary{positioning}
\usetikzlibrary{calc}
\usetikzlibrary{trees}

\tikzset{	bblock/.style 	= {draw, minimum width=8em, minimum height=2em, text height=1.5ex,text depth=0.25ex},
	io/.style 		= {bblock, trapezium, trapezium left angle=70, trapezium right angle=110},
	op/.style 		= {bblock, rectangle}, 
	if/.style		= {bblock, diamond},
	terminator/.style		= {bblock, rounded rectangle},
	portal/.style 	={bblock,circle,minimum size=3em},
}

% Source - https://tex.stackexchange.com/a/23094
% Posted by egreg
% Retrieved 2026-02-16, License - CC BY-SA 3.0

\newcommand*{\fvtextcolor}[2]{\textcolor{#1}{#2}}
\newcommand*{\tc}[2]{\textcolor{#1}{#2}}


\begin{document}
	 
	\mode<presentation>
	{
		\input{template.tex}	
	}
	
	


	
\QRFRAME
	

	
\mode<presentation>
{	
	
			
	    
		
		\frame{\maketitle}
		
		\begin{frame}\frametitle{План лекции}			
			\tableofcontents
		\end{frame}

	
}



\mode<article>
{
	\begin{center}
		
		\Large
		
		\title
		
		\LectionTheme
	\end{center}
}
	
\section{Системы сборки программ}
	
	
	
\begin{frame}[fragile]\frametitle{Система сборки}



	\begin{block}{Система сборки --- }
		это программный инструмент или набор инструментов, который автоматизирует процесс преобразования исходного кода (а также связанных ресурсов, зависимостей и других файлов) в исполняемые программы или библиотеки.
	\end{block} \pause
	
	Сборка (build) программ для С++ --- это процесс, обычно включает в себя следующие ключевые этапы, управляемые системой сборки:
	
	\begin{itemize}
		\item \textbf{Препроцессинг}: обработка директив {\verb|#include|}, макросов (\verb|#define|) и условной компиляции (\verb|#if|, \verb|ifdef|, и т.д.).
		\pause\item \textbf{Компиляция}: перевод исходного кода (.cpp, .cxx файлов) в объектные файлы (.o, .obj), содержащие машинный код.
		\pause\item \textbf{Компоновка} (линковка): объединение всех объектных файлов и статических/динамических библиотек в единый исполняемый файл (.exe) или библиотеку(.dll, .lib, .a, .so).
	\end{itemize}
	
\end{frame}
	

\section{Cmake}

\begin{frame}\frametitle{Способы работы с Cmake}
	
	
	\TCT{0.5}	
	{
		\centering
		
		\textbf{GUI}
		
		\IG[1]{cmake_gui.jpg}
	}
	{
		\centering
		
		\textbf{CLI}
		
		\IG[1]{cmake_cli.jpg}
	}
	
	

	
\end{frame}


\begin{frame}[fragile, t]\frametitle{Соберём первую программу}
	
	\begin{onlyenv}<+>
		
	\textbf{1. Необходимые файлы }
		
	\vspace{1cm}	

	\begin{columns}[t]
		\edef\colA{0.5}
		\begin{column}{\colA\linewidth}
			%Колнка1
			Файл с исходным кодом программы:
			\begin{block}{\ttfamily main.c}
			\begin{cppcode}
				#include <stdio.h>
								
				int main()
				{
					printf("Hello from CMake!");				
					return 0;
				}
			\end{cppcode}
			\end{block}
		\end{column}
		\begin{column}{\fpeval{1-\colA}\linewidth}
			Файл с настройкой проекта:
			\begin{block}{\ttfamily CMakeLists.txt}
			\begin{cmakecode}
				cmake_minimum_required(VERSION 4.1)
				project(First_project)
				add_executable(pr main.c)
			\end{cmakecode}
			\end{block}
		\end{column}
	\end{columns}

	\end{onlyenv}
	\begin{onlyenv}<+>
		\textbf{2. Работаем с Cmake}
		

		
		
		\begin{columns}[onlytextwidth]
			\edef\colA{0.5}
			\begin{column}{\colA\linewidth}
				%Колнка1
				\begin{tikzpicture}[remember picture]
					\node[anchor=south west](a){\IG[0.95]{cmake1}};
%					\draw[thick] (a.south west) grid (a.north east);
%					\draw[thin, opacity=0.5, step=0.5] (a.south west) grid (a.north east);
					
					\draw [red, thick] (1.6,5.25) -- (4, 5.25) node[below, midway] {\scriptsize диретория с \ttfamily CMakeLists.txt};
					
					\coordinate (P1) at (0.75,2.6);
					\coordinate (P4) at (2.5,2.3);
					
					
				\end{tikzpicture}
			\end{column}
			\begin{column}{\fpeval{1-\colA}\linewidth}
				%Колнка2
				\begin{tikzpicture}[remember picture]
					\node[anchor=south west](a){\IG[0.95]{cmake2}};
%					\draw[thick] (a.south west) grid (a.north east);
%					\draw[thin, opacity=0.5, step=0.5] (a.south west) grid (a.north east);
					
					\coordinate (P2) at (0.4,3.55);
					\coordinate (P3) at (4.5,0.4);
				\end{tikzpicture}
			\end{column}
			\begin{tikzpicture}[remember picture,overlay]
				\draw[red, -> , very thick] (P1) ..controls ++(30:2.5).. (P2);
			\end{tikzpicture}
			\begin{tikzpicture}[remember picture,overlay]
				\draw[red, -> , very thick] (P3) ..controls ++(180:3).. (P4);
			\end{tikzpicture}
		\end{columns}
		
	\end{onlyenv}
	
	\begin{onlyenv}<+>
		\textbf{3. Результат}
		
		\centering
		\IG[0.75]{cmake3}
		
		
	\end{onlyenv}

	
\end{frame}


\begin{frame}[fragile]\frametitle{Структура папок}

	\begin{columns}[]
		\edef\colA{0.45}
		\begin{column}{\colA\framewidth}
			%Колнка1
			\begin{Verbatim}[bgcolor=black,tabsize=4,formatcom=\color{white},gobble=4,fontsize=\small, commandchars=\\\{\}, xleftmargin=0em, frame=single,framerule=0pt ,rulecolor=\color{red},framesep=1em,fillcolor=\color{black}]
				
				Sample program
				¦   CMakeLists.txt \color{red}{<==}
				¦   main.c    \color{green}{<==}
				¦   
				L---build
					¦   ALL_BUILD.vcxproj
					¦   ALL_BUILD.vcxproj.filters
					¦   ALL_BUILD.vcxproj.user
					¦   CMakeCache.txt \color{cyan}{<==}
					¦   ...
					...   
					¦
					+---Debug
					¦       pr.exe   \color{yellow}{<==}
					¦       pr.pdb
					...
				
			\end{Verbatim}
		\end{column}
		\begin{column}{\fpeval{1-\colA}\framewidth}
			%Колнка2
			\textbf{Проверим программу:}
			
			\begin{Verbatim}[bgcolor=black,tabsize=4,formatcom=\color{white},gobble=4,fontsize=\footnotesize, xleftmargin=0em, frame=single,framerule=0pt ,rulecolor=\color{red},framesep=1em,fillcolor=\color{black}]
				
				PowerShell 7.6.0-preview.6
				
				PS Simple program> .\build\Debug\pr.exe
				Hello from CMake!
				
				PS Simple program>
				
			\end{Verbatim}
		\end{column}
	\end{columns}


	
	
\end{frame}

\begin{frame}[fragile]\frametitle{Теперь то же самое в консоли (Windows)}
		
		\begin{Verbatim}[bgcolor=black,tabsize=4,formatcom=\color{white},gobble=2,fontsize=\footnotesize, xleftmargin=0em, frame=single,framerule=0pt ,rulecolor=\color{red},framesep=1em,fillcolor=\color{black},commandchars=+\{\}]
		
		PS C:\Users\gavre\Desktop\Simple program> +color{orange}rm -r build -Force
		PS C:\Users\gavre\Desktop\Simple program> +color{orange} cmake -B build
		-- Building for: Visual Studio 18 2026
		-- Selecting Windows SDK version 10.0.26100.0 to target Windows 10.0.26200.
		-- The C compiler identification is MSVC 19.50.35724.0
		-- The CXX compiler identification is MSVC 19.50.35724.0
		...
		-- Build files have been written to: C:/Users/gavre/Desktop/Simple program/build 
		
		+pause{}PS C:\Users\gavre\Desktop\Simple program> +color{orange} cmake --build build
		...
		pr.vcxproj -> C:\Users\gavre\Desktop\Simple program\build\Debug\pr.exe
		
		+pause{}PS C:\Users\gavre\Desktop\Simple program> .\build\Debug\pr.exe
		Hello from CMake!  +color{red}{<==}

		
	\end{Verbatim}
	
\end{frame}


\begin{frame}[fragile]\frametitle{И еще раз (Linux)}
	
	\begin{Verbatim}[bgcolor=black,tabsize=4,formatcom=\color{white},gobble=2,fontsize=\footnotesize, xleftmargin=0em, frame=single,framerule=0pt ,rulecolor=\color{red},framesep=1em,fillcolor=\color{black},commandchars=+\{\}]
		
		gavreg@GvrNoutiL:.../Simple program$ +color{orange} cmake -B build_linux
		-- The C compiler identification is GNU 14.2.0
		-- The CXX compiler identification is GNU 14.2.0
		...
		-- Configuring done (6.4s)
		-- Generating done (0.3s)
		-- Build files have been written to: .../Simple program/build_linux
		+pause{}gavreg@GvrNoutiL:.../Simple program$ +color{orange} cmake --build build_linux
		...
		[100%] Built target pr
		gavreg@GvrNoutiL:.../Simple program$ +color{orange} tree build_linux/
		build_linux/
			├── CMakeCache.txt
			├── Makefile
			├── pr
		...
		+pause{}gavreg@GvrNoutiL:.../Simple program$ +color{orange}./build_linux/pr
		Hello from CMake!   +color{red}{<==}
		
	\end{Verbatim}
	

	
\end{frame}


\begin{frame}[fragile]\frametitle{Простые переменные}
	
	\begin{minted}[autogobble]{cmake}
		set(<variable> <value>... [PARENT_SCOPE])
	\end{minted}
	
	\verb|<variable>| --- имя переменной
	
	\verb|<value>| --- значение переменной 
	
	\verb|PARENT_SCOPE| --- родительская область видимости
	
	\begin{columns}[]
		\edef\colA{0.5}
		\begin{column}{\colA\linewidth}
			%Колнка1
			\begin{block}{CMakeLists.txt}
				\begin{cmakecode}
					cmake_minimum_required(VERSION 4.2)
					project("First project")
					
					set(A 1)
					message("${A}")
					set(A 1 2 3)
					message("${A}")
					
					add_executable(pr main.c)
				\end{cmakecode}
			\end{block}
		\end{column}
		\begin{column}{\fpeval{1-\colA}\linewidth}
			%Колнка2
			\begin{Verbatim}[bgcolor=black,tabsize=5,formatcom=\color{white},gobble=4,fontsize=\footnotesize, xleftmargin=0em, frame=single,framerule=0pt ,rulecolor=\color{red},framesep=1em,fillcolor=\color{black}, commandchars=+\{\}]
				
				PS ...> +color{orange} cmake -B build
				...
				1
				1;2;3
				... 
				
			\end{Verbatim}
		\end{column}
	\end{columns}
	
\end{frame}


\begin{frame}[fragile]\frametitle{Простые переменные}
	\begin{columns}[]
		\edef\colA{0.5}
		\begin{column}{\colA\linewidth}
			%Колнка1
			\begin{block}{CMakeLists.txt}
				\begin{cmakecode}
					cmake_minimum_required(VERSION 4.1)
					project("First project")
					
					set(X "ФГАОУ")
					set(Y "ВО МГТУ")
					set(Z "СТАНКИН")
					
					message(${X} ${Y} ${Z})
					message("${X} ${Y} ${Z}")
					
					add_executable(pr main.c)
				\end{cmakecode}
			\end{block}

		\end{column}
		\begin{column}{\fpeval{1-\colA}\linewidth}
			%Колнка2
				\begin{Verbatim}[bgcolor=black,tabsize=4,formatcom=\color{white},gobble=4,fontsize=\footnotesize, xleftmargin=0em, frame=single,framerule=0pt ,rulecolor=\color{red},framesep=1em,fillcolor=\color{black},commandchars=+\{\}]
				
				PS ...> +color{orange}  cmake -B build
				...
				ФГАОУВО МГТУСТАНКИН
				ФГАОУ ВО МГТУ СТАНКИН
				...
				-- Build files...
				
			\end{Verbatim}
		\end{column}
	\end{columns}	
\end{frame}

\begin{frame}[fragile]\frametitle{Кешированные переменные}
	
	\begin{minted}[autogobble]{cmake}
		#c версии 4.2
		set(CACHE{<variable>} [TYPE <type>] [HELP <helpstring>...] 
				[FORCE] VALUE [<value>...])
		
		#до версии 4.2
		set(<variable> <value>... CACHE <type> <docstring> [FORCE])
	\end{minted}
	
	\vspace{1em}
	
	\begin{tabular}{|l|l|}
		\hline
		\textbf{Тип} &  \textbf{Описание} \\\hline
		\Verb|BOOL|  &  Булево ON/OFF значение \\\hline
		\Verb|FILEPATH|  &  Путь к файлу  диске  \\\hline
		\Verb|PATH|  & Путь к папке на диске  \\\hline
		\Verb|STRING|  &  Строка\\\hline
		\Verb|INTERNAL|  & Строка для внутреннего использовани  \\\hline
	\end{tabular}

\end{frame}


\begin{frame}[fragile]\frametitle{Установка кешированных переменных через GUI}
	
	\begin{columns}[]
		\edef\colA{0.45}
		\begin{column}{\colA\linewidth}
			%Колнка1
			\begin{block}{CMakeLists.txt}
			\begin{cmakecode}
				cmake_minimum_required(VERSION 4.2.3)
				project("First project")
				
				set (CACHE{A} TYPE BOOL 
					FORCE VALUE ON)
				set (CACHE{B}  TYPE FILEPATH 
					FORCE VALUE "C:/Users/gavre")
				set (CACHE{C}  TYPE STRING 
					FORCE VALUE "lalal")
				
				add_executable(pr main.c)
			\end{cmakecode}
			\end{block}
		\end{column}
		\begin{column}{\fpeval{1-\colA}\linewidth}
			%Колнка2
			\IG{cmake_cache}
		\end{column}
	\end{columns}
		
\end{frame}

\begin{frame}[fragile]\frametitle{Кешированные переменные. Пример}
	\begin{columns}[]
		\edef\colA{0.5}
		\begin{column}{\colA\linewidth}
			%Колнка1
			\begin{block}{CMakeLists.txt}
				\begin{cmakecode}
					cmake_minimum_required(VERSION 4.1)
					project("First project")
					
					set(X "ФГАОУ" CACHE STRING  "Описание1")
					set(Y "ВО МГТУ" CACHE STRING "Описание2")
					set(Z "СТАНКИН" CACHE STRING "Описание3")
					message(${X} ${Y} ${Z})
					message("${X} ${Y} ${Z}")
					
					add_executable(pr main.c)
				\end{cmakecode}
			\end{block}
			
		\end{column}
		\begin{column}{\fpeval{1-\colA}\linewidth}
			%Колнка2
			\begin{Verbatim}[bgcolor=black,tabsize=4,formatcom=\color{white},gobble=4,fontsize=\footnotesize, xleftmargin=0em, frame=single,framerule=0pt ,rulecolor=\color{red},framesep=1em,fillcolor=\color{black},commandchars=+\{\}]
				
				PS ...> +color{orange}  cmake -B build
				...
				ФГАОУВО МГТУСТАНКИН
				ФГАОУ ВО МГТУ СТАНКИН
				...
				-- Build files... 
				
				+pause PS ...> +color{orange} cat .\build\CMakeCache.txt
				...
				//Описание1
				X:STRING=ФГАОУ
				
				//Описание2
				Y:STRING=ВО МГТУ
				
				//Описание3
				Z:STRING=СТАНКИН
				...
				
			\end{Verbatim}
		\end{column}
	\end{columns}	
\end{frame}

\begin{frame}[fragile]\frametitle{Изменение кешированных переменных}
	\begin{columns}[]
		\edef\colA{0.5}
		\begin{column}{\colA\linewidth}
			%Колнка1
			\begin{block}{CMakeLists.txt}
				\begin{cmakecode}
					cmake_minimum_required(VERSION 4.1)
					project("First project")
					
					message(${X} ${Y} ${Z})
					
					add_executable(pr main.c)
				\end{cmakecode}
			\end{block}
			
		\end{column}
		\begin{column}{\fpeval{1-\colA}\linewidth}
			%Колнка2
			\begin{Verbatim}[bgcolor=black,tabsize=4,formatcom=\color{white},gobble=4,fontsize=\footnotesize, xleftmargin=0em, frame=single,framerule=0pt ,rulecolor=\color{red},framesep=1em,fillcolor=\color{black},commandchars=+\{\}]
				
				PS ...> +color{orange}cmake -B build
				...
				ФГАОУ ВО МГТУ СТАНКИН
				...
				
			\end{Verbatim}
		\end{column}
	\end{columns}	
	
	\pause\vspace{0.3em}\rule{\linewidth}{0.4pt}
	
	\begin{columns}[]
		\edef\colA{0.5}
		\begin{column}{\colA\linewidth}
			%Колнка1
			\begin{block}{CMakeLists.txt}
				\begin{cmakecode}
					cmake_minimum_required(VERSION 4.1)
					project("First project")
					
					set(Z "им. Баумана" CACHE STRING "Фууу")
					message("${X} ${Y} ${Z}")
					
					add_executable(pr main.c)
				\end{cmakecode}
			\end{block}
			
		\end{column}
		
		\begin{column}{\fpeval{1-\colA}\linewidth}
			%Колнка2
			\begin{Verbatim}[bgcolor=black,tabsize=4,formatcom=\color{white},gobble=4,fontsize=\footnotesize, xleftmargin=0em, frame=single,framerule=0pt ,rulecolor=\color{red},framesep=1em,fillcolor=\color{black},commandchars=+\{\}]
				
				PS ...> +color{orange}cmake -B build
				...
				ФГАОУ ВО МГТУ СТАНКИН
				...+pause
				PS ...> +color{orange}rm -r build -Force
				PS ...> +color{orange}cmake -B build
				...
				им. Баумана
				...
				
			\end{Verbatim}
		\end{column}
	\end{columns}	
\end{frame}

\begin{frame}[fragile]\frametitle{Изменение кешированных переменных. Force}
	\begin{columns}[]
		\edef\colA{0.5}
		\begin{column}{\colA\linewidth}
			%Колнка1
			\begin{block}{CMakeLists.txt}
				\begin{cmakecode}
					...					
					set (CACHE{A} VALUE "Адын!")
					message("${A}")
					set (CACHE{A} VALUE "Дуа!")
					message("${A}")
					set (CACHE{A} VALUE "Тури!")
					message("${A}")
					...
				\end{cmakecode}
			\end{block}
			
		\end{column}
		\begin{column}{\fpeval{1-\colA}\linewidth}
			%Колнка2
			\begin{Verbatim}[bgcolor=black,tabsize=4,formatcom=\color{white},gobble=4,fontsize=\footnotesize, xleftmargin=0em, frame=single,framerule=0pt ,rulecolor=\color{red},framesep=1em,fillcolor=\color{black},commandchars=+\{\}]
				
				PS ...> +color{orange}rm -r build -Force
				PS ...> +color{orange}cmake -B build
				...
				Адын!
				Адын!
				Адын!
				...
				
			\end{Verbatim}
		\end{column}
	\end{columns}	
	
	\pause\vspace{0.3em}\rule{\linewidth}{0.4pt}
	
	\begin{columns}[]
		\edef\colA{0.5}
		\begin{column}{\colA\linewidth}
			%Колнка1
			\begin{block}{CMakeLists.txt}
				\begin{cmakecode}
					...
					set (CACHE{A} FORCE VALUE "Адын!")
					message("${A}")
					set (CACHE{A} FORCE VALUE "Дуа!")
					message("${A}")
					set (CACHE{A} FORCE VALUE "Тури!")
					message("${A}")
					...
				\end{cmakecode}
			\end{block}
			
		\end{column}
		
		\begin{column}{\fpeval{1-\colA}\linewidth}
			%Колнка2
			\begin{Verbatim}[bgcolor=black,tabsize=4,formatcom=\color{white},gobble=4,fontsize=\footnotesize, xleftmargin=0em, frame=single,framerule=0pt ,rulecolor=\color{red},framesep=1em,fillcolor=\color{black},commandchars=+\{\}]
				
				PS ...> +color{orange}cmake -B build
				...
				Адын!
				Дуа!
				Тури!
				...
				
			\end{Verbatim}
		\end{column}
	\end{columns}	
\end{frame}


\begin{frame}[fragile]\frametitle{Переменные окружения}
	
	\begin{minted}[autogobble]{cmake}
		set(ENV{<variable>} [<value>])
	\end{minted}
	
	\begin{columns}[]
		\edef\colA{0.5}
		\begin{column}{\colA\linewidth}
			%Колнка1
			\begin{block}{CMakeLists.txt}
			\begin{cmakecode}
				cmake_minimum_required(VERSION 4.2.3)
				project("First project")				
				
				set  (ENV{A} 1213)				
				message ("$ENV{A}")		
						
				execute_process(COMMAND cmd /C "echo %A%")
								
				message("$ENV{PATH}")
				
				add_executable(pr main.c)
			\end{cmakecode}
			\end{block}
		\end{column}
		\begin{column}{\fpeval{1-\colA}\linewidth}
			%Колнка2
			\begin{Verbatim}[bgcolor=black,tabsize=4,formatcom=\color{white},gobble=4,fontsize=\footnotesize, xleftmargin=0em, frame=single,framerule=0pt ,rulecolor=\color{red},framesep=1em,fillcolor=\color{black},commandchars=+\{\}]
				
				PS ...> +color{orange}cmake -B build
				...
				Моя переменная = 1213
				1213
				C:\Program Files\WindowsApps\Microsoft.P
				owerShell_7.5.4.0_x64__8wekyb3d8bbwe;C:\
				ProgramFiles\ImageMagick-7.1.1-Q16-HDRI;
				...
				
			\end{Verbatim}
		\end{column}
	\end{columns}
	
\end{frame}

\begin{frame}[fragile]\frametitle{Основные встроенные переменные CMake}
	\small
	\begin{tabular}{>{\ttfamily}l l l}
		\hline
		\textbf{Переменная} & \textbf{Описание} & \textbf{Пример} \\
		\hline
		CMAKE\_SOURCE\_DIR & Корень исходников & /home/user/proj \\
		CMAKE\_BINARY\_DIR & Корень сборки & /home/user/proj/build \\
		CMAKE\_CURRENT\_SOURCE\_DIR & Текущий исходник & /home/user/proj/src \\
		CMAKE\_VERSION & Версия CMake & 3.27.1 \\
		CMAKE\_SYSTEM & ОС & Linux, Windows \\
		CMAKE\_SYSTEM\_NAME & Имя ОС & Linux \\
		CMAKE\_SYSTEM\_PROCESSOR & Архитектура & x86\_64 \\
		CMAKE\_BUILD\_TYPE & Тип сборки & Debug \\
		CMAKE\_INSTALL\_PREFIX & Путь установки & /usr/local \\
		PROJECT\_NAME & Имя проекта & MyProject \\
		PROJECT\_SOURCE\_DIR & Корень проекта & /home/user/proj \\
		EXECUTABLE\_OUTPUT\_PATH & Путь для bin & /home/user/proj/build/bin \\
		CMAKE\_C\_STANDARD & Стандарт C & 90, 99, 11, 17, 23 \\
		CMAKE\_CXX\_STANDARD & Стандарт C++ & 98, 11, 14, 17, 20, 23 \\
		CMAKE\_C\_STANDARD\_REQUIRED & Треб. стандарт C & TRUE / FALSE \\
		CMAKE\_CXX\_STANDARD\_REQUIRED & Треб. стандарт C++ & TRUE / FALSE \\
		\hline
	\end{tabular}
	
	\footnotesize{\textit{Полный список:} \ttfamily cmake --help-variable-list}
	
	\footnotesize{\textit{Информация о переменной:} \ttfamily cmake --help-variable ИМЯ\_ПЕРЕМЕННОЙ}
\end{frame}

\begin{frame}[fragile]\frametitle{Встроенные переменные связанные с компилятором}
	\small
	\begin{tabular}{>{\ttfamily}l l l}
		\hline
		\textbf{Переменная} & \textbf{Описание} & \textbf{Пример} \\
		\hline
		CMAKE\_C\_COMPILER & Компилятор C & /usr/bin/gcc \\
		CMAKE\_CXX\_COMPILER & Компилятор C++ & /usr/bin/g++ \\
		CMAKE\_COMPILER\_IS\_GNU\_CC & GNU C компилятор? & TRUE / FALSE \\
		CMAKE\_COMPILER\_IS\_GNU\_CXX & GNU C++ компилятор? & TRUE / FALSE \\
		CMAKE\_C\_COMPILER\_ID & Идент-ор. C компилятора & GNU, Clang, MSVC \\
		CMAKE\_CXX\_COMPILER\_ID & Идент-ор. C++ компилятора & GNU, Clang, MSVC \\
		CMAKE\_C\_COMPILER\_VERSION & Версия C компилятора & 11.4.0 \\
		CMAKE\_CXX\_COMPILER\_VERSION & Версия C++ компилятора & 11.4.0 \\
		CMAKE\_C\_FLAGS & Флаги для C компилятора & -O2 -Wall \\
		CMAKE\_CXX\_FLAGS & Флаги для C++ компилятора & -O2 -Wall -std=c++11 \\
		CMAKE\_C\_FLAGS\_DEBUG & Флаги для Debug (C) & -g \\
		CMAKE\_CXX\_FLAGS\_DEBUG & Флаги для Debug (C++) & -g \\
		CMAKE\_C\_FLAGS\_RELEASE & Флаги для Release (C) & -O3 -DNDEBUG \\
		CMAKE\_CXX\_FLAGS\_RELEASE & Флаги для Release (C++) & -O3 -DNDEBUG \\
		CMAKE\_EXE\_LINKER\_FLAGS & Флаги линковщика & -static -lpthread \\
		CMAKE\_LINKER & Линковщик & /usr/bin/ld \\
		\hline
	\end{tabular}
	
	
\end{frame}

\begin{frame}[fragile]\frametitle{Условия}
	
	\begin{cmakecode}
		if(<условие>)
			Команды
		elseif(<другое условие>) # необязательно. можно повторить
			Команды
		else()              # необязательно
			Команды
		endif()
	\end{cmakecode}
	
	\hfill
	
	\mintinline{cmake}|if(<константа>)|
	
	Истина --- 1 ON YES Y TRUE $\neq 0$
	
	Ложь --- 0 OFF NO FALSE N IGNORE NOTFOUND \"\" "...-NOTFOUND"
	
	\hfill
	
	\mintinline{cmake}|if(<перееменная>)|
	
	Истина --- существующая переменная, не являющаяся ложность константой
	
	Ложь --- во всех остальных случаях
	
	\hfill
	
	\mintinline{cmake}|if("<строка>")|
	
	Истина --- Если является истинной константой
	
	Ложь --- во всех остальных случаях
	
\end{frame}


\begin{frame}\frametitle{Логические операторы}
	
	\mintinline{cmake}{if(NOT <условие>)} --- отрицание
	
	\hfill
	
	\mintinline{cmake}{if(<условие1> AND <условие2>)} --- логическое И
	
	\hfill
	
	\mintinline{cmake}{if(<условие1> OR <условие2>)} --- логическое ИЛИ
	
	\hfill
	
	\mintinline{cmake}{if((условие1) AND (условие2 OR (услвие3)))} --- сложыне условия
	
\end{frame}

\begin{frame}[fragile]\frametitle{Компараторы}
	

	\mintinline{cmake}{if(<variable|string> ХХХ <variable|string>)} 
	
	\hfill
	
	\begin{tabularx}{\textwidth}{>{\ttfamily}rcX}
		XXX: &&\\
		LESS & --- & строго меньше \\
		GREATER & --- & строго больше \\
		EQUAL & --- & равно \\
		LESS\_EQUAL & --- & меньше или равно (3.7+) \\
		GREATER\_EQUAL & --- & больше или равно (3.7+) \\
		STRLESS & --- & лексикографически меньше \\
		STRGREATER & --- & лексикографически больше \\
		STREQUAL & --- & лексикографически равно \\
		STRLESS\_EQUAL & --- & лексикографически меньше или равно (3.7+) \\
		STRGREATER\_EQUAL & --- & лексикографически больше или равно (3.7+) \\
		MATCHES & --- & соответствует регулярному выражению (группы~() сохраняются в CMAKE\_MATCH\_<n> (2.6+)) \\
	\end{tabularx}


	
	
\end{frame}

\begin{frame}[fragile]\frametitle{Сравнение версий и путей}
	
	\mintinline{cmake}{if(<variable|string> XXX <variable|string>)} 
	
	\hfill
	
	\textbf{Сравнение версий:}
	\begin{tabularx}{\textwidth}{>{\ttfamily}rcX}
		VERSION\_LESS & --- & строго меньше \\
		VERSION\_GREATER & --- & строго больше \\
		VERSION\_EQUAL & --- & равно \\
		VERSION\_LESS\_EQUAL & --- & меньше или равно (3.7+) \\
		VERSION\_GREATER\_EQUAL & --- & больше или равно (3.7+) \\
	\end{tabularx}
	
	\smallskip
	Формат версии: major[.minor[.patch[.tweak]]], пропущенные компоненты~=~0
	
	\hfill
	
	\textbf{Сравнение путей (3.24+):}
	\begin{tabularx}{\textwidth}{>{\ttfamily}rcX}
		PATH\_EQUAL & --- & покомпонентное сравнение путей \\
	\end{tabularx}
	
	\smallskip
	Коллапсирование разделителей:
	\mintinline{cmake}{if("/a//b/c" PATH_EQUAL "/a/b/c")} = true
	
\end{frame}


\begin{frame}[fragile]\frametitle{Циклы. Foreach}
	
	\begin{minted}[autogobble]{cmake}
		foreach(<loop_var> <items>)
		<commands>
		endforeach()
	\end{minted}
	
	
	\begin{columns}[]
		\edef\colA{0.5}
		\begin{column}{\colA\linewidth}
			%Колнка1
			\begin{block}{CMakeLists.txt}
				\begin{cmakecode}
					foreach(A "Раз два три")
					message("Итерация = ${A}")
					endforeach()
					
					foreach(A Раз два три)
					message("Итерация = ${A}")
					endforeach()
					
					set(X Вася Петя Маша)
					
					foreach(A ${X})
					message("Имя = ${A}")
					endforeach()
				\end{cmakecode}
			\end{block}
		\end{column}
		\begin{column}{\fpeval{1-\colA}\linewidth}
			%Колнка2
			\begin{Verbatim}[bgcolor=black,tabsize=4,formatcom=\color{white},gobble=4,fontsize=\footnotesize, xleftmargin=0em, frame=single,framerule=0pt ,rulecolor=\color{red},framesep=1em,fillcolor=\color{black},commandchars=+\{\}]
				
				PS ...> +color{orange}cmake -B build
				...
				Итерация = Раз два три
				Итерация = Раз
				Итерация = два
				Итерация = три
				Имя = Вася
				Имя = Петя
				Имя = Маша
				...
				
			\end{Verbatim}
		\end{column}
	\end{columns}
	
\end{frame}

	
\end{document}
